\DocumentMetadata{
  pdfversion=2.0,
  pdfstandard=ua-2,
  lang=en-US,
  tagging=on,
  tagging-setup={math/setup=mathml-SE,tabsorder=structure}
}
\documentclass[11pt,letterpaper]{article}
\usepackage[margin=0.68in,includefoot]{geometry}
\usepackage{newcomputermodern}
\usepackage{amsmath,amscd,mathtools}
\usepackage{tikz,adjustbox}
\providecommand{\square}{\mdlgwhtsquare}
\usepackage[hidelinks]{hyperref}
\hypersetup{
  pdftitle={Topology PhD Exam, August 2022},
  pdfauthor={University of Florida Department of Mathematics},
  pdfsubject={Graduate mathematics examination},
  pdfdisplaydoctitle=true,
  bookmarksopen=true
}
\setcounter{secnumdepth}{0}
\setlength{\parindent}{0pt}
\setlength{\parskip}{0.28em}
\setlength{\emergencystretch}{3em}
\setlength{\leftmargini}{2.15em}
\setlength{\leftmarginii}{2.65em}
\setlength{\leftmarginiii}{3em}
\pagestyle{empty}
\raggedbottom
\allowdisplaybreaks
\NewTaggingSocketPlug{block/list/label}{examproblem}{%
  \tagstructbegin{tag=Lbl}\tagstructend%
  \tagstructbegin{tag=\LBody}%
  \tagstructbegin{tag=\ProblemHeadingTag,title-o={\ProblemHeadingText}}%
  \tagmcbegin{tag=\ProblemHeadingTag,actualtext={\ProblemHeadingText}}%
  #2%
  \tagmcend%
  \tagstructend%
}
\newenvironment{examproblems}[2]{%
  \def\ProblemHeadingTag{#2}%
  \AssignTaggingSocketPlug{block/list/label}{examproblem}%
  \begin{enumerate}%
  \setcounter{enumi}{#1}%
  \renewcommand{\labelenumi}{\textbf{\theenumi.}}%
  \setlength{\topsep}{0.35em}%
  \setlength{\partopsep}{0pt}%
  \setlength{\itemsep}{0.48em}%
  \setlength{\parsep}{0.18em}%
}{\end{enumerate}}
\newenvironment{examsubparts}{%
  \AssignTaggingSocketPlug{block/list/label}{default}%
  \begin{enumerate}%
  \setlength{\topsep}{0.18em}%
  \setlength{\partopsep}{0pt}%
  \setlength{\itemsep}{0.16em}%
  \setlength{\parsep}{0.1em}%
}{\end{enumerate}}
\newenvironment{examinstructions}{%
  \par\begingroup\itshape
}{%
  \par\endgroup\vspace{0.18em}
}
\makeatletter
\renewcommand\subsection{\@startsection{subsection}{2}{0pt}%
  {-1.25ex plus -.25ex minus -.1ex}%
  {0.55ex plus .1ex}%
  {\normalfont\large\bfseries}}
\renewcommand{\@maketitle}{%
  \newpage\null\vskip -1em%
  {\footnotesize\bfseries
    \href{https://www.ufl.edu/}{University of Florida}, \href{https://math.ufl.edu/}{Department of Mathematics}\par}%
  \vskip -0.5em%
  \tagstructbegin{tag=H1,title={Topology PhD Exam, August 2022}}%
  \tagpdfparaOff%
  \tagmcbegin{tag=H1}%
    {\LARGE\bfseries\href{https://gma.math.ufl.edu/exams/topology-phd/}{Topology PhD Exam}, August 2022\par}%
  \tagmcend%
  \tagpdfparaOn%
  \tagstructend%
  \vskip 0.25em%
  \tagmcbegin{artifact}%
    \hrule height 1pt%
  \tagmcend%
  \vskip 1em%
}
\makeatother
\title{Topology PhD Exam, August 2022}
\author{}
\date{}
\begin{document}
\maketitle
\thispagestyle{empty}
\begin{examinstructions}
Work the following problems and show all work. Support all statements to the best of your ability.
\end{examinstructions}
\begin{examproblems}{0}{H2}
\def\ProblemHeadingText{Problem 1.}
\item
Suppose~\(M\) is a compact 5-manifold with \(H_0(M;\mathbb{Z})=\mathbb{Z}\), \(H_1(M;\mathbb{Z})=\mathbb{Z}_3\), and \(H_2(M;\mathbb{Z})=\mathbb{Z}\). Is~\(M\) orientable? What are \(H_3(M;\mathbb{Z})\), \(H_4(M;\mathbb{Z})\), and \(H_5(M;\mathbb{Z})\)?
\def\ProblemHeadingText{Problem 2.}
\item
Let~\(M\) be a compact, connected, nonorientable 3-manifold. Prove that \(H_1(M;\mathbb{Z})\) is infinite.
\def\ProblemHeadingText{Problem 3.}
\item
Let~\(X\) be a connected metric space with metric~\(d\) and let \(p\in X\). Prove that if \(X\setminus\{p\}\ne\varnothing\), then \(X\setminus\{p\}\) is not compact.
\def\ProblemHeadingText{Problem 4.}
\item
Let \(\Sigma_g\) denote the closed surface of genus~\(g\), and let \(\Sigma_{g,r}\) denote the surface of genus~\(g\) with~\(r\) boundary components (i.e., \(\Sigma_g\) with~\(r\) discs removed). Which of the surfaces \(\Sigma_{g,r}\) can admit a map \(f:\Sigma_{g,r}\to\Sigma_{g,r}\) homotopic to the identity that does not have a fixed point?
\def\ProblemHeadingText{Problem 5.}
\item
Let~\(X\) be the space obtained by filling in two discs in the torus, as shown (the discs lie inside the surface of the torus). Compute \(H_\bullet(X;\mathbb{Z})\).
\par\smallskip
\begin{center}
\begin{adjustbox}{max width=.8\linewidth,max totalheight=6\baselineskip}
\begin{tikzpicture}[
  alt={Diagram of a torus shown as a flattened oval with a small central opening. Two black-outlined, diagonally hatched disks fill openings on the left and right sides of the torus.},
  line cap=round,
  line join=round,
  x=1cm,
  y=1cm
]
% Left filling disc
\fill[black!20]
    (-4.4,0)
    .. controls (-3.65,0.85) and (-2.10,1.25) .. (-1.05,0.15)
    .. controls (-1.75,-0.75) and (-3.45,-0.80) .. cycle;

\begin{scope}
    \clip
        (-4.4,0)
        .. controls (-3.65,0.85) and (-2.10,1.25) .. (-1.05,0.15)
        .. controls (-1.75,-0.75) and (-3.45,-0.80) .. cycle;
    \foreach \x in {-4.7,-4.3,-3.9,-3.5,-3.1,-2.7,-2.3,-1.9,-1.5,-1.1}
        \draw[line width=0.45pt]
            (\x,-1.05) -- ++(0.62,2.15);
\end{scope}

\draw[line width=1pt]
    (-4.4,0)
    .. controls (-3.65,0.85) and (-2.10,1.25) .. (-1.05,0.15)
    .. controls (-1.75,-0.75) and (-3.45,-0.80) .. cycle;

% Right filling disc
\fill[black!20]
    (1.15,0.15)
    .. controls (2.15,1.05) and (3.75,1.15) .. (4.4,0)
    .. controls (3.95,-0.85) and (2.15,-0.80) .. cycle;

\begin{scope}
    \clip
        (1.15,0.15)
        .. controls (2.15,1.05) and (3.75,1.15) .. (4.4,0)
        .. controls (3.95,-0.85) and (2.15,-0.80) .. cycle;
    \foreach \x in {0.9,1.3,1.7,2.1,2.5,2.9,3.3,3.7,4.1}
        \draw[line width=0.45pt]
            (\x,-1.05) -- ++(0.62,2.15);
\end{scope}

\draw[line width=1pt]
    (1.15,0.15)
    .. controls (2.15,1.05) and (3.75,1.15) .. (4.4,0)
    .. controls (3.95,-0.85) and (2.15,-0.80) .. cycle;

% Outer boundary of the torus
\draw[line width=1.2pt]
    (-4.4,0)
    .. controls (-4.45,1.65) and (-2.25,2.45) .. (0,2.40)
    .. controls (2.40,2.38) and (4.20,1.65) .. (4.4,0);

\draw[line width=1.2pt]
    (4.4,0)
    .. controls (4.15,-1.10) and (2.25,-1.45) .. (0,-1.35)
    .. controls (-2.45,-1.48) and (-4.25,-1.05) .. (-4.4,0);

% Central hole
\draw[line width=1.15pt]
    (-1.05,0.15)
    .. controls (-0.55,0.72) and (0.65,0.72) .. (1.15,0.15)
    .. controls (0.62,-0.48) and (-0.55,-0.48) .. cycle;
\end{tikzpicture}
\end{adjustbox}
\end{center}
\smallskip
\end{examproblems}
\begin{examinstructions}
Answer the following with complete definitions, statements, or short proofs.
\end{examinstructions}
\begin{examproblems}{5}{H2}
\def\ProblemHeadingText{Problem 6.}
\item
Prove or give a counterexample: If \(A\subset X\) is path-connected, then the closure \(\overline{A}\) is path-connected.
\def\ProblemHeadingText{Problem 7.}
\item
Compute \(\chi(\mathbb{R}P^2\times\mathbb{C}P^3\times S^4)\).
\def\ProblemHeadingText{Problem 8.}
\item
Give the definition of a normal topological space. Show that a smooth manifold is a normal space.
\def\ProblemHeadingText{Problem 9.}
\item
Prove that a continuous bijection from a compact space to a Hausdorff space is a homeomorphism.
\def\ProblemHeadingText{Problem 10.}
\item
Prove that if \(m\ne n\), then \(\mathbb{R}^m\) is not homeomorphic to \(\mathbb{R}^n\).
\def\ProblemHeadingText{Problem 11.}
\item
Prove that the torus \(T^2=S^1\times S^1\) is not homotopy equivalent to \(S^1\vee S^1\vee S^2\).
\def\ProblemHeadingText{Problem 12.}
\item
State the Tietze Extension Theorem.
\def\ProblemHeadingText{Problem 13.}
\item
Describe all the connected covering spaces \(E\to\mathbb{R}P^2\vee\mathbb{R}P^3\).
\def\ProblemHeadingText{Problem 14.}
\item
Does the following exact sequence of abelian groups necessarily split? Prove or give a counterexample.
\[
0\to\mathbb{Z}_2\to A\to\mathbb{Z}_2\to 0
\]
\def\ProblemHeadingText{Problem 15.}
\item
Compute the integral homology of the space \(\mathbb{C}P^2\times\mathbb{R}P^2\).
\end{examproblems}
\end{document}
